\section{Bidirectional early charging under an explicit settlement rule}
\label{sec:v2g}
The two-batch model also permits an exact extension in which assigned buses
can discharge before the second batch. This changes both the physical
opportunity set and the tariff interpretation. We assume ideal efficiency
one, symmetric charging and discharging acceptance, and the \emph{same}
posted price for signed net purchases and exports in a market interval.
An asymmetric buy/sell tariff or a lossy battery is a different model.

Retain $n$ A services in $[0,1]$, $n$ B services in $[2,3]$, capacity
$B>0$, consumptions $0<a_i,b_j\le B$, and the fixed allowed successor
graph $\Gamma$ of Section~\ref{sec:heterogeneous}. Every singleton remains
available; allowed transitions have zero time, energy and cost. A matching
of size $m$ uses exactly $N_m=2n-m$ assigned buses, each with fee $f>0$,
full initial inventory and full terminal replenishment. We first derive
the projection with service-assigned buses, then allow identical paid idle
buses below. In particular, a B-only bus is assigned to service,
starts full and is available at the depot throughout the early window.

All assigned buses can use the single early connector during $[1,2]$ in
either direction at power magnitude $K\ge0$. Its unit-window resource is
\emph{gross} throughput: if bus $v$ has signed net early purchase $q_v$,
then $\sum_v|q_v|\le K$ must be realizable. A constraint on $|\sum_vq_v|$
alone would incorrectly allow charging and discharging to cancel connector
time. Terminal replenishment in $[3,4]$ uses one charging-only connector
of power $H\ge0$. Each bus accepts those corresponding powers. There
is no other charging, reserve, auxiliary draw, taper or switching time.

Net early load $E=\sum_vq_v$ may now be negative. The nonnegative fleet-load
convention in Section~\ref{sec:framework} is therefore replaced explicitly
for this section: $F$ must be defined and finite on all admitted signed
loads $(E,S-E)$. Alternatively, declare fixed background demand and an
incremental cost $F(e)=C(d+e)-C(d)$ whose domain contains the whole feasible
signed-load set. Export settlement still uses the same linear rate as
purchases. There is no implicit extension of a cost defined only on
nonnegative fleet demand. If aggregate export is prohibited, impose $E\ge0$
explicitly; the lower thresholds below are then additionally clipped at zero.

Write $S=\sum_i a_i+\sum_j b_j$, $P=\sum_i a_i$, and let $\nu$
be the maximum matching cardinality of $\Gamma$. Retain the matching values
\[
 r_m=\min_{\substack{M\subseteq\Gamma\\|M|=m}}
       \sum_{(i,j)\in M}[a_i+b_j-B]_+,
 \qquad 0\le m\le\nu,
\]
with $r_0=0$. The matching machinery is unchanged; the available assigned
battery inventory and gross connector budget create the new thresholds.

\begin{lemma}[Minimum gross energy for a fixed matching]\label{lem:v2ggross}
For a fixed allowed matching $M$ of size $m$, put
\[
 R=\sum_{(i,j)\in M}[a_i+b_j-B]_+,\qquad T=S-N_mB.
\]
Before the shared connector constraint, signed early charges have total
$E\in[T,P]$. At a prescribed such $E$, their minimum gross throughput is
\begin{equation}\label{eq:v2ggross}
 \min_{\sum_vq_v=E}\sum_v|q_v|=\max\{E,2R-E\}.
\end{equation}
The minimum is attained by one monotone action of the chosen sign per bus.
\end{lemma}

\begin{theorem}[Signed-load projection and joint frontier]\label{thm:v2g}
Under the above assumptions, define
\begin{equation}\label{eq:v2gthresholds}
 t_m=\max\{S-H,\ S-(2n-m)B,\ 2r_m-K\},\qquad
 U=\min\{P,K\},\qquad c_m=f(2n-m).
\end{equation}
The complete physical projection is exactly the union of
\begin{equation}\label{eq:v2gprojection}
 \{(E,S-E,c_m):t_m\le E\le U\}
\end{equation}
over $0\le m\le\nu$ with $t_m\le U$. The model is feasible exactly
when $t_0\le U$. The sequence $t_m$ is nondecreasing and discretely convex.

When feasible, let $m_{\max}=\max\{m\le\nu:t_m\le U\}$ and
$m_0=\max\{m\le m_{\max}:t_m=t_0\}$. The hull's lowest operating-cost
boundary starts at $(t_0,c_{m_0})$, visits $(t_m,c_m)$ for
$m_0<m\le m_{\max}$, and is horizontal at $c_{m_{\max}}$ on
$[t_{m_{\max}},U]$. Physical thresholds on collinear hull edges remain
part of this statement.
\end{theorem}

\begin{corollary}[Support and physical attainment]\label{cor:v2gsupport}
Suppose $t_0\le U$ and $G(E)=F(E,S-E)$ is finite, continuous and strictly
convex on $[t_0,U]$. Let $E^*$ be the unique minimizer of the hull
operating boundary plus $G$. Then, for $f>0$,
\[
 D=\CH\quad\Longleftrightarrow\quad
 E^*\in\{t_0\}\cup\{t_m:m_0<m\le m_{\max}\}
             \cup[t_{m_{\max}},U].
\]
Admissible own-load support additionally requires the framework's convex-
supply, dual-attainment and price-set qualification on the declared signed
domain. At $f=0$, the empty-matching branch has zero operating cost over
the entire feasible interval and the gap is zero, so the positive-fee
strict-interior criterion must not be applied.
\end{corollary}

\subsection{Why the projection is exact}
\begin{proof}[Proof of Lemma~\ref{lem:v2ggross}]
A paired bus has signed charge in $[a_i+b_j-B,a_i]$: early headroom
gives the upper bound, and enough inventory for B gives the lower bound.
An A-only bus has $q\in[a_i-B,a_i]$; a B-only bus has
$q\in[b_j-B,0]$. These endpoint conditions keep inventory in $[0,B]$
before and after the remaining service. A monotone action between those
inventories is also feasible throughout the early window. Terminal demand
is that bus's service consumption minus $q$, between zero and $B$.

All upper endpoints are nonnegative and sum to $P$. The sum of positive
lower endpoints is $R$, while the sum of all lower endpoints is
$T=S-N_mB$. Thus total possible negative excursion is $R-T$. For
$E\ge R$, set negative coordinates to zero and distribute purchases
among the independent nonnegative intervals, whose sum spans $[R,P]$.
Gross energy is then $E$. For $T\le E<R$, keep mandatory positive
coordinates at their lower endpoints and distribute discharge of magnitude
$R-E$ among the remaining negative intervals. This is possible exactly
because $R-E\le R-T$ and has gross energy $2R-E$.

Conversely, total positive purchases are at least $R$, so any allocation
with net $E$ has gross energy at least $2R-E$ and at least $E$.
The constructions attain both regimes in \eqref{eq:v2ggross}. Additional
charge/discharge cycles cannot reduce gross energy and are unnecessary.
\end{proof}

\begin{proof}[Proof of Theorem~\ref{thm:v2g}]
The lemma and early gross budget give $T\le E\le P$, $E\le K$
and $E\ge2R-K$. All buses are present for the full early window;
their monotone actions serialize in at most one hour when $K>0$.
For $K=0$, the inequalities force $R=0$, $E=0$ and zero gross energy,
without division. Every bus's terminal refill is nonnegative and at most
$B$. They serialize at $H$ exactly when $S-E\le H$; at $H=0$ no
positive terminal demand can be met. With a nonempty positive B batch,
$E\le P<S$ makes that zero-terminal-power case infeasible.

For a fixed matching, the exact interval is consequently
\[
 [\max\{S-H,S-N_mB,2R(M)-K\},U].
\]
Only $R(M)$ varies at fixed cardinality. Its attained minimum $r_m$
gives the widest interval; all such intervals have common upper endpoint
$U$. Their union is \eqref{eq:v2gprojection}. Since $r_m$ is
nondecreasing, its threshold minimum is $t_0$; feasibility is precisely
$t_0\le U$, attained by a standalone fleet when needed.

Standard allowed-edge bipartite matching convexity gives discretely convex
$r_m$. Interpolate it linearly on $[0,\nu]$. Taking the maximum of
the resulting convex function $2r-K$, the affine function $S-(2n-m)B$
and the constant $S-H$ gives a convex function whose integer values are
$t_m$. Each component is nondecreasing, so the thresholds are also
nondecreasing. Zero increments can occur only in an initial block.
After retaining its cheapest cardinality, the positive threshold increments
are nondecreasing. Frontier slopes $-f/(t_m-t_{m-1})$ are therefore
nondecreasing, followed by zero. Each endpoint is physical and each edge
is its whole-plan mixture; every physical branch lies above this convex
graph. This proves the lower-boundary assertion without introducing
nonexistent or resource-infeasible cardinalities.
\end{proof}

\begin{proof}[Proof of Corollary~\ref{cor:v2gsupport}]
Strict convexity of $G$ gives one minimizing hull load. At a retained
threshold or on the final horizontal segment, a physical plan attains the
frontier. In an adjacent open threshold interval, $f>0$ makes its
interpolated cost strictly below every physical branch at that load.
The compact finite union of physical intervals then has a strictly larger
minimum. The classical support criterion applies only with the stated
signed-domain and price qualifications. At fee zero the empty matching
physically attains the entire interval at zero operating cost.
\end{proof}

\subsection{Identical paid idle buses}
An idle bus has the same capacity, initial full inventory, terminal-full
requirement, connector access and fee as a service bus. Its early signed
charge lies in $[-B,0]$ and its terminal refill is $-q$. It has no separate
energy source or export tariff.

\begin{proposition}[Idle buses under a total fleet cap]
Allow identical paid idle buses, with total fleet at most $2n$. The complete
cost-labelled projection in Theorem~\ref{thm:v2g} is unchanged.
\end{proposition}
\begin{proof}
Consider a matching of size $m$ with $k$ idle buses. Its total fleet is
$N=2n-m+k$, so the cap implies $k\le m$. The box argument of
Lemma~\ref{lem:v2ggross} still applies: idle buses add only negative
headroom. The matching's interval is therefore
\[
 [\max\{S-H,S-NB,2R-K\},U],
\]
at operating cost $fN$. Delete any $k$ matching edges and use the resulting
$m-k$ matching without idle buses. This fleet has the same size $N$ and
fee. Deleting edges cannot increase $R$, while the inventory term $S-NB$
and upper endpoint $U$ stay fixed. Its feasible interval contains the
original interval. Thus every projected point with idle buses has an
identical load and cost without them. The reverse inclusion is immediate.
\end{proof}

For the two examples below, $H=S$. Even if the fleet cap is removed,
fleets larger than $2n$ cannot improve the lower frontier: terminal
capacity forces $E\ge0$, and the $2n$ standalone buses realize every
$E\in[0,U]$ at fee $2nf$. Any larger fleet has the same upper endpoint
$U$ and a higher fee when $f>0$. It is therefore dominated at each
feasible load. The literal projection acquires higher-cost layers, but
its lower frontier and support conclusions are unchanged. This argument
does not apply when $H>S$: additional idle inventory can then enable
negative net early loads and may be useful.

\subsection{Persistence and removal of an intermediate charge requirement}
For the earlier anchor with A energies $(5,15)$, B energies $(14,15)$,
battery capacity 20, $K=15$ and $H=49$, the matching values are
$(r_0,r_1,r_2)=(0,0,9)$ and the new thresholds remain
\[
 (t_0,t_1,t_2)=(0,0,9).
\]
Although individual buses gain discharging options, the entire cost-labelled
projection is unchanged. At the earlier fee and supply objective, its
previous positive gap therefore survives analytically. The terminal bound
here enforces $E\ge0$, so this particular conclusion does not require
evaluating the old supply function at negative fleet loads. This is a proof
consequence, not a new V2G implementation qualification or experiment.

Bidirectional charging can also remove a positive gap. Take complete
crossmatching, A energies $(1,1,19)$, B energies $(2,10,19)$, capacity
20, $K=2$, $H=52$ and fee $f=1$. Here $S=52$ and
$(r_0,r_1,r_2,r_3)=(0,0,0,1)$. In the charge-only model, the cheapest
physical operating cost is 4 on $0\le E<1$ and 3 on $1\le E\le2$.
Its hull boundary is $4-E$ on $[0,1]$, followed by the constant 3.
For the convex supply cost
\[
 F(E,L)=E^2+\frac{E+L}{5},\qquad
 G(E)=E^2+\frac{52}{5},
\]
the physical minima at $E=0$ and $E=1$ tie, whereas the hull minimum
is at $E=1/2$. Consequently,
\[
 D=\frac{72}{5},\qquad \CH=\frac{283}{20},\qquad
 D-\CH=\frac14.
\]

With bidirectional charging, all four thresholds $t_m$ are zero, so the
three-bus operating cost 3 is physical throughout $[0,2]$. In particular,
pair the services as $(1,19)$, $(1,10)$ and $(19,2)$ and assign signed
early charges $0,-1,1$, respectively. Their net load is zero and gross
throughput is 2. Their terminal refills are $20,12,20$, totaling 52;
all intermediate inventories respect capacity. Thus
\[
 D=\CH=3+\frac{52}{5}=\frac{67}{5},\qquad E^*=0.
\]
At the own-load price $\nabla F(0,52)=(1/5,1/5)$, every feasible plan
pays $52/5$ for net energy, and the three-bus plan minimizes operating
cost. Hence it is supported under the stated settlement rule. Discharge
from another assigned bus offsets a mandatory positive early purchase,
using the gross connector budget in both directions. These exact values
follow analytically from the projection theorem; they are not a numerical
implementation qualification.

Losses cannot be incorporated by rescaling the early-energy threshold alone.
Let $\eta_c,\eta_d$ denote charge and discharge efficiencies, $X$ early
grid exports and $E$ early net grid purchase. Full replenishment requires
terminal grid purchase
\[
 L=\frac{S}{\eta_c}-E+
       \left(\frac{1}{\eta_c\eta_d}-1\right)X.
\]
For strict round-trip loss, the same net early load can therefore have
different terminal demand. The lossless one-dimensional projection does
not transfer directly. Under the common-window assumptions and nonnegative
prices or marginal supply cost, within-bus cycles can be removed without
worsening the objective, but cross-bus redistribution cannot generally be
removed in that way. Section~\ref{sec:lossy} derives a joint efficient-frontier
reduction with common losses; the lossless scalar support classification
does not follow from that reduction.

The result concerns lossless common-window physics and the specified net
settlement. It neither establishes novelty of V2G or matching theory nor
covers degradation, asymmetric prices, realistic travel, or the public
37-service case. That case remains charge-only and unchanged. No numerical
qualification of this bidirectional extension is claimed.
